\documentclass[11pt,a4paper]{appletmanual}
\usepackage[spanish,es-noquoting,es-nodecimaldot]{babel}
\manuallang{es}

\appletname{Funciones de Demanda}
\appletsub{De dónde salen las curvas de demanda. Mueve un precio o la renta, y
cada canasta óptima deja una marca: el rastro que se va formando es la función
de demanda.}
\appletdir{demand-functions}

\begin{document}
\manualtitlepage{Manual del applet · Microeconomía I · Universidad de los Andes\\
Traducido de un applet de GeoGebra.\\
Código MIT. Se abre en el navegador; no hay nada que instalar.}

\manualtoc
\thispagestyle{empty}
\clearpage
\setcounter{page}{1}

% ==========================================================================
\section{Qué es esto}

Una curva de demanda no es un dato. Es el resultado de resolver el problema del
consumidor una vez por cada precio, y de apuntar la respuesta.

En un curso eso se dice, se dibuja la curva ya hecha, y se pasa al siguiente
tema. Este applet hace lo otro: pone el problema del consumidor a la izquierda,
un eje vacío a la derecha, y deja que el estudiante mueva el precio. Cada
canasta óptima deja una marca en el panel derecho. Al cabo de un rato hay una
curva ahí, y el estudiante la ha construido punto por punto.

\figher[1.0]{revealed}{Después de barrer $p_x$. Los puntos morados del panel
izquierdo son la senda precio-consumo; los del derecho son la curva de demanda
que se acaba de construir, con la curva exacta revelada por debajo.}

\begin{amkey}
El rastro y la curva no son dos cosas: la curva revelada pasa exactamente por
los puntos registrados. Ver que coinciden es la diferencia entre saber la
definición de una función de demanda y entenderla.
\end{amkey}

\subsection{Para quién}

Para el tema de demanda, después de haber visto el óptimo del consumidor. Es el
puente entre los dos. Sirve además para bienes normales, inferiores y Giffen,
que suelen quedar como una nota al pie y aquí se pueden fabricar y mirar.

%% Secciones comunes a todos los manuales, edición española.
%% Se incluyen con \input{common-es}. Lo único que cambia entre applets es
%% \appletfolder, que es también el nombre de la carpeta y de la ruta en el sitio.

\section{Cómo abrirlo}

No hay que instalar nada. La aplicación entera es un solo archivo
\code{index.html}: no descarga librerías, no llama a ningún servidor y no
guarda nada fuera del navegador.

\subsection{Las tres maneras}

\begin{description}
\item[Doble clic sobre el archivo.] Es lo más rápido. Busque
  \code{\appletfolder/index.html} y ábralo. El navegador lo muestra desde el disco,
  con la barra de direcciones empezando por \code{file://}. Funciona sin
  conexión desde el primer momento.
\item[Desde la web.] Si alguien ya lo publicó, la dirección termina en
  \code{/\appletfolder/}. Una vez cargada la página, puede desconectarse: no
  necesita la red para nada más.
\item[Servido desde una carpeta.] Para una clase, poner la carpeta en
  cualquier servidor estático basta. Con Python instalado:
  \begin{quote}\code{python3 -m http.server 8000}\end{quote}
  y luego \code{http://localhost:8000/\appletfolder/} en el navegador.
\end{description}

\begin{amnote}
Guardar la página con \key{Ctrl}\,+\,\key{S} (o \key{Cmd}\,+\,\key{S} en
Mac) deja una copia propia que sigue funcionando sin conexión y que se puede
llevar en una memoria USB o repartir por correo.
\end{amnote}

\subsection{Idioma y tema}

Arriba a la derecha hay dos pares de botones.

\begin{ctrltable}{Control & Qué hace}
ES / EN & Cambia el idioma de toda la interfaz al vuelo. Nada se recalcula:
          puede cambiarlo en mitad de una explicación sin perder el estado. \\
Oscuro / Claro & Cambia el tema. El oscuro es el diseño original; el claro
          existe para proyectores y para imprimir. \\
\end{ctrltable}

Las dos elecciones se recuerdan en el navegador, así que la próxima vez se abre
como la dejó.

\subsection{Qué navegador}

Cualquiera que sea razonablemente reciente: Chrome, Edge, Firefox o Safari, en
ordenador, tableta o teléfono. En pantallas estrechas los paneles se apilan uno
sobre otro en lugar de ponerse al lado. No hace falta cuenta, ni permisos, ni
extensiones.


% ==========================================================================
\section{La pantalla}

\fig[1.0]{interface}{La pantalla completa. Izquierda: el problema del
consumidor. Derecha: la demanda que se va construyendo.}

\begin{description}
\item[Intervalos.] El rango que se considera para cada variable. Es a la vez el
  recorrido del deslizador y el tramo que dibuja el eje horizontal del panel
  derecho. Es el control que más se olvida y el que más problemas resuelve.
\item[Preferencias.] Cinco familias, la última escribible.
\item[Parámetros.] Los de la familia elegida.
\item[Variable que varía.] Cuál de $p_x$, $p_y$ o $m$ se mueve. Las otras dos
  quedan fijas, y los puntos registrados pertenecen solo a esa.
\item[Precios y renta.] $p_x$, $p_y$, $m$.
\item[Descubrir la demanda.] Los botones que registran, barren y limpian.
\item[Capas.] Qué se dibuja.
\end{description}

\subsection{Los dos paneles}

\begin{description}
\item[Elección · plano $(x, y)$.] La recta presupuestaria, la curva de
  indiferencia por la canasta óptima, la canasta misma, y la senda
  precio-consumo que se va formando.
\item[Demanda.] La cantidad óptima frente a la variable que se mueve. Su
  cabecera dice cuál es. Cuando la variable es la renta, el panel se llama
  \ui{Curva de Engel}.
\end{description}

El conmutador \ui{Un panel / Dos} de la cabecera derecha separa la demanda de
$x$ de la de $y$ en dos gráficos, lo que hace falta cuando las dos tienen
escalas muy distintas.

\fig[1.0]{engel}{Con la renta como variable que varía, el panel derecho pasa a
ser la curva de Engel.}

\subsection{El tema oscuro}

\fig[0.95]{dark}{Los dos paneles en el tema oscuro.}

% ==========================================================================
\section{Las cinco familias}

Cada una trae una demanda en forma cerrada, salvo la personalizada, que se
resuelve numéricamente.

\begin{ctrltable}{Familia & Utilidad y demanda}
Cobb--Douglas & $U = x^{\alpha} y^{1-\alpha}$, con
  $x^\ast = \alpha m/p_x$ e $y^\ast = (1-\alpha) m/p_y$. La demanda de cada
  bien no depende del precio del otro: es el caso limpio. \\
CES & $U = (x^{r} + y^{r})^{1/r}$. Con $r$ cerca de 1 son casi sustitutos
  perfectos; con $r$ muy negativo, casi complementarios. \\
Lineal (sustitutos) & $U = a x + y$. La demanda es de esquina: todo en $x$ si
  $a > p_x/p_y$, todo en $y$ si es al revés. \\
X inferior (Giffen) & $U = \ln(x - \underline{x}) -
  \tfrac{1+b}{b}\ln(\bar y - y)$ con $b = (1-\beta)/\beta$. Es la familia con
  la que se fabrica el caso Giffen. \\
Personalizada & La $U(x, y)$ que escriba. \\
\end{ctrltable}

\begin{ctrltable}{Parámetro & De qué familia}
$\alpha$ & Cobb--Douglas: la participación de $x$ en el gasto. \\
$r$ & CES: el exponente de sustitución. \\
$a$ & Lineal: cuánto vale $x$ frente a $y$. \\
$\beta$, $\underline{x}$, $\bar y$ & X inferior: el peso, el mínimo de
  subsistencia de $x$, y el techo de $y$. \\
\end{ctrltable}

% ==========================================================================
\section{Descubrir la demanda}

Es el corazón del applet. Cuatro controles:

\begin{ctrltable}{Control & Qué hace}
Marcar al mover & Registra un punto automáticamente cada vez que se mueve la
  variable elegida. Es lo más cómodo para una demostración en clase. \\
Registrar punto & Guarda la canasta actual. Un punto, deliberadamente. \\
Barrido & Recorre la variable de un extremo a otro del intervalo y va
  marcando. Pulsar otra vez lo para. \\
Limpiar & Borra los puntos de la variable actual. \\
\end{ctrltable}

Los puntos pertenecen a la variable que se estaba moviendo. Cambiar de $p_x$ a
$m$ no mezcla los rastros: cada uno guarda el suyo.

\fig[1.0]{swept}{A mitad de un barrido de $p_x$: los puntos ya registrados y la
canasta óptima actual.}

\subsection{Las capas}

\begin{ctrltable}{Capa & Qué dibuja}
Puntos registrados & Las marcas que se han ido dejando. \\
Unir los puntos & Los conecta con segmentos. \\
Revelar la curva & Dibuja la función de demanda exacta, calculada, por debajo
  de los puntos. Es la comprobación. \\
Senda precio-consumo & En el plano $(x, y)$, el rastro de las canastas
  óptimas. \\
Recta presupuestaria & La restricción actual. \\
Curva de indiferencia & La que pasa por la canasta óptima. \\
Demanda de x / Demanda de y & Cuál de las dos demandas se dibuja en el panel
  derecho. \\
Resaltar tramo Giffen & Marca en morado los tramos con
  $\partial x^\ast/\partial p_x > 0$. \\
Retícula & La cuadrícula de fondo. \\
\end{ctrltable}

% ==========================================================================
\section{Intervalos y ejes}

\begin{amwatch}
Los deslizadores de $p_x$, $p_y$ y $m$ están \emph{acotados por el intervalo}
que se haya fijado en el grupo \ui{Intervalos}. Si intenta bajar $p_y$ por
debajo de su mínimo de intervalo no se mueve. Baje primero el mínimo del
intervalo. Es la causa más frecuente de que un deslizador «no responda».
\end{amwatch}

El eje vertical del panel derecho tiene dos modos:

\begin{description}
\item[Eje vertical fijo] (por defecto). El tope es $m_{\text{máx}}/p_{\text{mín}}$,
  la cantidad máxima comprable con los intervalos vigentes. La ventaja es que el
  eje no se mueve al arrastrar, así que dos barridos se pueden comparar a ojo.
\item[Automático.] El eje se ajusta a la curva. Hace falta cuando la cantidad de
  interés es pequeña frente a ese tope —el caso Giffen, por ejemplo, donde
  $x^\ast$ ronda $0{,}3$ y el tope está en $8$.
\end{description}

% ==========================================================================
\section{Normal, inferior, Giffen}

Las tres lecturas de la derecha lo resuelven todo:

\begin{ctrltable}{Lectura & Qué dice}
$\partial x^\ast/\partial p_x$ & La pendiente de la demanda. Negativa es lo
  normal; positiva es Giffen. \\
$\partial x^\ast/\partial m$ & El efecto renta. Negativo significa bien
  inferior. \\
Tipo & \emph{normal} o \emph{inferior}, según el signo del anterior. \\
\end{ctrltable}

Y el veredicto de arriba nombra el caso:

\begin{description}
\item[x es un bien normal aquí.] La demanda baja cuando sube su precio y sube
  con la renta.
\item[x es un bien inferior aquí.] $\partial x^\ast/\partial m < 0$. Ser
  inferior es \emph{necesario} para ser Giffen, pero no basta.
\item[x es un bien Giffen aquí.] $\partial x^\ast/\partial p_x > 0$. El efecto
  renta domina al de sustitución, cosa que solo puede pasar si $x$ es inferior.
\end{description}

\fig[1.0]{giffen}{El caso Giffen. La demanda sube con el precio, y todo el
tramo está resaltado. Observe el panel izquierdo: la senda precio-consumo sube
casi en vertical, porque el consumidor apenas puede reducir $x$ por debajo de
su mínimo de subsistencia.}

\fig[0.85]{giffen-verdict}{El veredicto en el caso Giffen.}

% ==========================================================================
\section{Actividades para clase}

\begin{activity}{Construir una demanda a mano}
\amset Valores de partida (Cobb--Douglas). \ui{Variable que varía} en $p_x$.
Encienda \ui{Marcar al mover}.
\amexp Con $p_x = 1{,}83$: $x^\ast = 1{,}68$, $y^\ast = 2{,}11$, gasto
$5{,}4$ de $5{,}4$. Mueva $p_x$ despacio de un extremo al otro. Cada canasta
deja una marca a la derecha.
\par\smallskip
Cuando tenga quince o veinte puntos, encienda \ui{Revelar la curva}. La curva
calculada pasa exactamente por ellos. Pregunte a la clase qué es entonces una
función de demanda: no un dato, sino el registro de veinte problemas de
optimización resueltos.
\end{activity}

\begin{activity}{El gasto no se mueve}
\amset Cobb--Douglas, barrido de $p_x$.
\amexp La lectura \ui{Gasto} dice siempre $5{,}4 / 5{,}4$, y $y^\ast$ no cambia
al mover $p_x$. Con Cobb--Douglas cada bien se lleva una fracción fija de la
renta: $x^\ast = \alpha m / p_x$ no depende de $p_y$, e $y^\ast$ no depende de
$p_x$.
\par\smallskip
Mueva $\alpha$ y vea las dos participaciones cambiar. Es la propiedad que hace
de Cobb--Douglas el caso de examen, y conviene que se vea antes de darla por
supuesta.
\end{activity}

\begin{activity}{La curva de Engel}
\amset \ui{Variable que varía} en $m$. Barra.
\amexp El panel derecho pasa a llamarse \ui{Curva de Engel}. Con
Cobb--Douglas sale una recta por el origen: la demanda es proporcional a la
renta.
\par\smallskip
Cambie a \ui{X inferior} y repita. La curva de Engel se dobla hacia abajo:
más renta, menos $x$. Ahí está la definición de bien inferior, dibujada.
\end{activity}

\begin{activity}{Inferior no es Giffen}
\amset Familia \ui{X inferior}, valores de partida.
\amexp $x^\ast = 1{,}08$, $y^\ast = 3{,}12$,
$\partial x^\ast/\partial p_x = -0{,}422$ y
$\partial x^\ast/\partial m = -0{,}068$. El tipo es \emph{inferior} y el
veredicto lo dice, pero la demanda sigue teniendo pendiente negativa.
\par\smallskip
Pregunte si eso contradice algo. No: ser inferior es necesario para ser Giffen,
no suficiente. Falta que el efecto renta sea lo bastante fuerte, y para eso hay
que empobrecer al consumidor en términos reales.
\end{activity}

\begin{activity}{Fabricar un bien Giffen}
\amset Familia \ui{X inferior}. En \ui{Intervalos}, baje el mínimo de $p_y$ a
$0{,}2$. Después ponga $p_y = 0{,}35$. Quite el \ui{Eje vertical fijo}.
\amexp Ahora $x^\ast = 0{,}30$, $y^\ast = 10{,}54$, y
$\partial x^\ast/\partial p_x = +3{,}1\times 10^{-4}$: \emph{positiva}. El
veredicto dice «x es un bien Giffen aquí» y el tramo sale resaltado.
\par\smallskip
Barra $p_x$ y mire la curva subir. Lo que ha hecho es abaratar tanto el otro
bien que $x$ se ha convertido en el bien de subsistencia del consumidor: cuando
$x$ sube de precio, se empobrece, y al empobrecerse compra más de lo barato.
Es la historia de la patata irlandesa, y hace falta este ajuste para verla.
\end{activity}

\begin{activity}{Los dos casos extremos}
\amset Familia \ui{Lineal (sustitutos)}. Barra $p_x$.
\amexp La demanda es un escalón: todo en $x$ hasta que $p_x/p_y$ cruza $a$, y
entonces cero de golpe. No hay tramo intermedio.
\par\smallskip
Cambie a \ui{CES} y baje $r$ hacia $-6$: las curvas se suavizan hasta parecerse
a las de complementarios. Los dos extremos acotan lo que puede hacer una
demanda, y el resto de las familias caen en medio.
\end{activity}

\begin{activity}{La senda precio-consumo}
\amset Cualquier familia. Encienda \ui{Senda precio-consumo} y barra $p_x$.
\amexp En el plano $(x, y)$ aparece el rastro de las canastas óptimas. Con
Cobb--Douglas es una horizontal, porque $y^\ast$ no depende de $p_x$. Con la
familia inferior sube casi en vertical.
\par\smallskip
La curva de demanda del panel derecho es esa misma senda, con $p_x$ en el eje
horizontal en vez de $x$. Son la misma información en dos sistemas de
coordenadas.
\end{activity}

\begin{activity}{Su propia función de utilidad}
\amset Familia \ui{Personalizada}. Escriba \code{x\^{}0.3*y\^{}0.7}, luego
\code{min(2x,y)}, luego \code{x+4ln(y)}.
\amexp La demanda se resuelve numéricamente y el rastro se construye igual. Con
\code{min(2x,y)} las dos demandas se mueven juntas; con la cuasilineal, la
demanda de $y$ deja de depender de la renta en cuanto esta pasa de un umbral.
\par\smallskip
Es el momento de pedir a la clase que prediga la forma antes de barrer.
\end{activity}

% ==========================================================================
\section{Sintaxis de expresiones}

\begin{ctrltable}{Elemento & Qué se admite}
Operadores & \code{+ - * / \^{} ( )} \\
Funciones & \code{sqrt ln log exp abs min max sin cos} \\
Constantes & \code{e}, \code{pi} \\
Multiplicación implícita & \code{2xy} es \code{2*x*y} \\
Variables & Solo \code{x} e \code{y} \\
\end{ctrltable}

% ==========================================================================
\section{Cómo se calcula}

Las cuatro familias con nombre traen su demanda en forma cerrada, deducida a
mano. La personalizada no puede, así que su óptimo se busca numéricamente sobre
la recta presupuestaria: un barrido denso y un refinamiento por sección áurea,
igual que en los otros applets de esta colección. Es robusto donde una
condición de primer orden no encuentra nada —esquinas y quiebres incluidos.

Las derivadas $\partial x^\ast/\partial p_x$ y $\partial x^\ast/\partial m$ que
aparecen en las lecturas son diferencias centradas sobre la propia función de
demanda, no derivadas simbólicas: así valen igual para las familias cerradas y
para la que escriba el usuario.

% ==========================================================================
\section{Problemas frecuentes}

\begin{ctrltable}{Síntoma & Qué pasa}
Un deslizador de precio no se mueve & Está topando con su intervalo. Cambie el
  mínimo o el máximo en \ui{Intervalos}. \\
La curva sale pegada al eje & El eje vertical está fijo en
  $m_{\text{máx}}/p_{\text{mín}}$ y la cantidad es pequeña. Quite \ui{Eje
  vertical fijo}. \\
Los puntos desaparecen al cambiar de variable & No desaparecen: cada variable
  guarda su propio rastro. Vuelva a la anterior y ahí están. \\
El panel derecho está ilegible & Las dos demandas tienen escalas muy distintas.
  Pase a \ui{Dos} paneles, o apague una de las dos capas. \\
No consigo el caso Giffen & Hacen falta las tres cosas: la familia
  \ui{X inferior}, un $p_y$ bajo (bajando antes su intervalo) y mirar el signo
  de $\partial x^\ast/\partial p_x$, no el de $\partial x^\ast/\partial m$. \\
\end{ctrltable}

%% Secciones comunes a todos los manuales, edición española.
%% Se incluyen con \input{common-es}. Lo único que cambia entre applets es
%% \appletfolder, que es también el nombre de la carpeta y de la ruta en el sitio.

\section{Publicarlo para una clase}

Sirve cualquier alojamiento estático: GitHub Pages, Netlify, un directorio web
de la universidad, o incluso una carpeta compartida servida por HTTP. Lo único
que hay que subir es el archivo \code{index.html} dentro de una carpeta con el
nombre que quiera que aparezca en la dirección.

El repositorio trae un flujo de trabajo de GitHub Actions
(\code{.github/workflows/pages.yml}) que publica todos los applets a la vez en
cada empuje a la rama principal. Para activarlo, en el repositorio:
\emph{Settings\,$\rightarrow$\,Pages\,$\rightarrow$\,Source: GitHub Actions}.
Es lo único que no puede hacer el propio flujo, porque crear un sitio de Pages
requiere permisos de administración que un \code{GITHUB\_TOKEN} nunca recibe.

\begin{amnote}
Para repartir el applet a estudiantes sin conexión fiable, mándeles el archivo
\code{index.html} directamente. Pesa menos que una fotografía y se abre con
doble clic.
\end{amnote}

\section{Licencia y procedencia}

El código de la aplicación es MIT. Puede usarlo, modificarlo y repartirlo, en
clase o fuera de ella, citando la procedencia.

Este applet es una reescritura autónoma de un applet original de GeoGebra. La
reescritura no es una traducción literal: donde el original tenía un error de
construcción, una división por cero alcanzable con los deslizadores, o un objeto
muerto heredado del applet del que se copió, esta versión lo corrige y lo
documenta. La sección \emph{Diferencias con el original de GeoGebra} enumera
esos cambios uno por uno.

Todo está escrito a mano y sin dependencias: el analizador de expresiones, la
derivación simbólica, el trazado de curvas de nivel y el dibujo. En particular
no se usa \code{eval}, de modo que lo que escribe un estudiante en una casilla
de texto no puede convertirse en código, y una política de seguridad de
contenidos estricta no rompe la página.


\end{document}
